\documentclass[11pt]{article}

\usepackage[a4paper,margin=2.2cm]{geometry}
\usepackage[T1]{fontenc}
\usepackage{lmodern}
\usepackage{microtype}
\usepackage{amsmath,amssymb,amsthm,mathtools}
\usepackage{booktabs}
\usepackage{array}
\usepackage{enumitem}
\usepackage{xcolor}
\usepackage{hyperref}

\hypersetup{
  colorlinks=true,
  linkcolor=blue!55!black,
  citecolor=green!45!black,
  urlcolor=blue!55!black,
  pdftitle={Coordination Is Not Simplicity}
}

\newtheorem{definition}{Definition}[section]
\newtheorem{theorem}[definition]{Theorem}
\newtheorem{proposition}[definition]{Proposition}
\newtheorem{example}[definition]{Example}
\newtheorem{counterexample}[definition]{Counterexample}
\theoremstyle{remark}
\newtheorem{remark}[definition]{Remark}

\newcommand{\MHC}{\mathsf{MHC}}
\newcommand{\Comp}{\mathsf{Completion}}
\newcommand{\Weak}{\mathsf{Weak}}
\newcommand{\Code}{\mathsf{Code}}
\newcommand{\Route}{\mathsf{Route}}
\newcommand{\Step}{\mathsf{Step}}
\newcommand{\Lean}{\textsc{Lean}}

\title{\textbf{Coordination Is Not Simplicity:}\\
Hierarchical Complexity, Weakness, Razors, and Governed Uptake}
\author{Oru\v{z}i \and Zarathustra Goertzel}
\date{Working paper, August 2026}

\begin{document}
\maketitle

\begin{abstract}
Hierarchical coordination, semantic freedom, weighted distinction, and
description length are often discussed as if they were interchangeable notions
of complexity.  We formalize them as separate observations of richer objects
and prove that none determines the others without additional premises.  Our
starting point is the Model of Hierarchical Complexity (MHC): an organization
is a chain when every permutation of its component actions has the same
outcome, and a coordination when some permutation changes the outcome.  We
extend this finitary account to arbitrary occurrence types and ordinal ranks,
recover the published natural-number theory as a specialization, and isolate
the hypotheses behind its finite cost law.

The formalization connects MHC to typed completion fibres, Bennett weakness,
Goertzel's quantale-valued weakness, code-relative selection, and proof-relevant
GSLT execution.  Independent completion factors as a product and independent
local execution commutes; relational coupling alone need not create an MHC
coordination.  A four-axis model gives bidirectional non-factorization witnesses
for hierarchical rank, completion cardinality, a nonuniform quantale score,
and code length.  Bennett's recoding result is then scoped precisely: it
obstructs a representation-invariant strict shorter-code selector, but does not
refute necessity-qualified ontological parsimony, reliability theorems for a
fixed model class, or fixed-vocabulary semantic invariance.

Finally, we formalize Crawford and Hammer's no-bypass condition for epistemic
resolve and test a proposed MHC reading of GSLT structure.  The latter does not
yield a context-free stage assignment.  One indexed naturality square is a
chain under endpoint observation and a coordination under proof-relevant route
observation.  Hierarchical order is therefore earned from an explicit outcome
observer, actual execution witnesses, and homogeneous child ranks; it is not a
label attached to a syntax, language transport, or software layer.  All stated
results are mechanized in Lean 4 with positive and negative controls.
\end{abstract}

\tableofcontents

\section{Four questions that should not share one scalar}

Suppose a system coordinates two procedures, admits many compatible futures,
has a short encoding, and receives a high score under a weighted valuation.
Four different questions have been asked:

\begin{enumerate}[label=(\arabic*)]
  \item Does permuting the procedures change their observed outcome?
  \item How many completions remain compatible with the system's commitments?
  \item Which distinctions matter under the declared valuation layer?
  \item How long is the system's description in the declared codebook?
\end{enumerate}

The first is operational, the second extensional, the third valuation-relative,
and the fourth representational.  A reduction between them can be useful, but
it must be a theorem with visible hypotheses.  Treating them as synonyms hides
exactly the distinctions needed for composition, generalization, and safe
self-modification.

Our method is therefore information-first.  We retain an action tree before
its ordinal rank, a completion fibre before its cardinality, an event before
its quantale value, a codebook before its length, and an update path before a
Boolean gate verdict.  Numerical readouts are derived only after their domain
and observer have been named.

\section{Infinite-primary hierarchical complexity}

\subsection{Schedules, chains, and coordinations}

Let $I$ be an arbitrary type of action occurrences.  A schedule is a
permutation $\sigma : I \simeq I$, and a schedule semantics is a map
\[
  s : (I\simeq I) \longrightarrow O
\]
to a declared outcome type $O$.

\begin{definition}[Chain and coordination]
A schedule semantics $s$ is a \emph{chain} when
\[
  \forall \sigma\,\tau,\quad s(\sigma)=s(\tau).
\]
It is a \emph{coordination} when
\[
  \exists \sigma\,\tau,\quad s(\sigma)\ne s(\tau).
\]
\end{definition}

This is the permutation criterion presented by Commons and Pekker
\cite{commonspekker2008}.  It is sharper than saying that a higher action
``non-arbitrarily organizes'' lower actions: an outcome and two schedules must
exhibit the organization or its invariance.

The ambient formalization permits arbitrary $I$.  Its actions form
well-founded, arbitrarily branching trees.  A simple action has rank zero.  An
invariant compound action takes the supremum of its child ranks; an
order-sensitive compound action takes the ordinal successor of that supremum:
\[
\begin{aligned}
 h(\mathsf{simple}) &= 0,\\
 h(\mathsf{chain}(a_i)_{i\in I}) &= \sup_{i\in I} h(a_i),\\
 h(\mathsf{coordination}(a_i)_{i\in I})
   &= \operatorname{succ}\!\left(\sup_{i\in I} h(a_i)\right).
\end{aligned}
\]

The last equation distinguishes successor coordination from limit formation.
A countably branching chain whose child ranks range over all finite ordinals
has rank $\omega$ without performing one additional coordination.

\subsection{Recovering the finitary theory}

Finite branching is a property of the ambient action tree, not the definition
of action.  Under finite branching, every rank lies below $\omega$ and hence is
a natural number.  A recursive admissibility predicate recovers the additional
Commons--Pekker condition that coordinated children have equal order.

The Lean development proves the familiar consequences inside this
specialization: order trichotomy and transitivity, maximum order for chains,
successor order for homogeneous coordinations, and a no-gaps predecessor
subaction.  Commons and Pekker's concrete examples are also computed rather
than merely named:

\begin{itemize}
  \item the separated operations $1+2$ and $3\times4$ are schedule invariant;
  \item the two schedules of $2\times(3+4)$ compute $14$ and $10$ and therefore
        form a coordination.
\end{itemize}

\subsection{What the exponential law does and does not say}

For a finite, atomically decomposed, homogeneous binary hierarchy, the minimum
number of simple leaves needed at order $n$ is
\[
  \varphi_n = 2^n.
\]
The formal result has two parts: every such action has at least $2^n$ simple
leaves, and a binary tower attains the bound.  Thus $2^n$ is a minimum theorem,
not a definition of order.

Two negative controls prevent an illicit transfinite extrapolation.  The
countably branching limit chain has $\aleph_0$ leaves, while
$2^{\aleph_0}>\aleph_0$.  Moreover, an inhabited atomless decomposition is not
well-founded and admits no strictly decreasing ordinal rank.  Consequently,
finite leaf counting is not an invariant of gunky or infinitely branching
organization.

\section{Completion geometry and level introduction}

\subsection{Principal fibres before Bennett cardinality}

Let $(P,\leq)$ be a preorder.  The principal completion fibre below $p$ is
\[
  \Comp(p) := \{q:P\mid q\leq p\}.
\]
This order-theoretic object is primary.  Bennett's finite weakness is recovered
as the cardinal readout $|\Comp(p)|$ when the fibre is finite.

For the product order on $P\times Q$, the formalization constructs an exact
equivalence
\[
  \Comp(p,q) \simeq \Comp(p)\times\Comp(q).
\]
Finite weakness therefore multiplies under independent composition.  An
interacting composition is represented instead by a compatibility-indexed
subfibre
\[
  \{(x,y):\Comp(p)\times\Comp(q)\mid C(x,y)\}.
\]
Strict compatibility can make this inclusion non-surjective.

The distinction between geometry and dynamics is essential.  A strict
relational coupling may still be schedule invariant.  Conversely, two
transformations on a shared state form a coordination exactly when their
composites disagree at the declared initial state.

\begin{theorem}[Binary operational bridge]
For a binary process with transformations $L,R$ and initial state $x$,
endpoint schedule semantics is a chain exactly when
\[
  R(L(x))=L(R(x)),
\]
and a coordination exactly when the two sides differ.
\end{theorem}

Component-local transformations on a product commute automatically.  With
equal-ranked children, the corresponding chain preserves their common rank;
a noncommuting coordination raises it by one.

\subsection{HC2 and C2 against task generations}

The existing Bennett task-refinement relation is transitive.  It therefore
cannot itself be a unit-ranked immediate-child relation.  A three-task example
has both a direct one-edge refinement witness and a mediated two-edge witness
between the same endpoints.  More generally, if two refinement edges compose,
then assigning a rank increase of exactly one to every broad refinement edge
produces inconsistent rank equations.

This is not a defect in Bennett's relation.  It identifies the additional data
needed for intrinsic levels.  A \emph{unit-ranked presentation} chooses an
immediate-step relation, proves that each step raises rank by one, and proves
that each chosen step is a valid Bennett refinement.  Exact paths in this
presentation have endpoint-determined length and a no-gaps spectrum.

The resulting verdicts are:

\begin{itemize}
  \item C2 holds for realized witness-path lengths, but these lengths are not
        intrinsic when the broad relation has shortcuts;
  \item C2 is intrinsic for a unit-ranked level presentation;
  \item HC2 is earned for equal-depth sources under one common target;
  \item systems terminology or broad refinement alone does not imply HC2.
\end{itemize}

\section{Four independent axes}

One finite candidate type retains four informative fields: an MHC action, a
completion set, a weighted pair event, and a codeword.  Four scalar readouts
are then defined:
\[
  h,\qquad |\Comp|,\qquad q,\qquad \ell.
\]
Here $h$ is hierarchical rank, $|\Comp|$ is Bennett weakness, $q$ is a declared
nonuniform quantale value, and $\ell$ is code length.

\begin{theorem}[Pairwise non-factorization]
For each of the six unordered pairs among
\[
  h,\qquad |\Comp|,\qquad q,\qquad \ell,
\]
there are
candidates with equal first readout and unequal second readout, and candidates
with equal second readout and unequal first readout.
\end{theorem}

All witnesses inhabit the same candidate space.  The codewords used in the
description-length witnesses are prefix-incomparable, so the separation is not
an artifact of one word extending another.

The negative result does not ban reductions.  It requires their premises to
be stated.  For the complete pair event $S\times S$,
\[
  |S\times S|=|S|^2,
\]
so pair cardinality factors through Bennett weakness and preserves and reflects
its finite order.  For the diagonal event $\{(s,s):s\in S\}$, cardinality is
exactly $|S|$.  Arbitrary nonuniform valuations need satisfy neither reduction.

\section{Ockham steelmanned and Bennett scoped}

\subsection{Several razors, not one target}

``Ockham's razor'' names a family of doctrines.  At least the following should
be separated:

\begin{itemize}
  \item necessity-qualified ontological economy;
  \item economy of explanatory assumptions or causes;
  \item evidentially constrained theory choice;
  \item a Bayesian prior favoring a declared simplicity measure;
  \item predictive-risk or class-capacity control justified by a reliability
        theorem;
  \item algorithmic coding criteria;
  \item MDL or MML model-plus-data message length.
\end{itemize}

Bennett's finite recoding argument targets a nontrivial strict
shorter-code-is-better selector when it is proposed as a
representation-invariant ideal \cite{bennett2026}.  It does not by itself
refute every item in the list.

The formal ontological profile first requires explanatory sufficiency and only
then minimizes unnecessary entities.  A theory with one necessary entity is
optimal over a theory with an additional redundant entity.  A theory with zero
entities but inadequate explanatory power is inadmissible.  Since this profile
does not inspect a codebook, it is invariant under recoding.

\subsection{Finite dyadic approximation and recoding}

For a nonempty finite hypothesis type and a strictly positive probability mass
$p$, the formalization constructs an explicit prefix codebook whose normalized
dyadic weights $q$ satisfy
\[
  \frac{1}{2}p(h) < q(h) \leq 2p(h)
\]
for every hypothesis $h$.  The construction uses equal-length injective tags
and minimally rounded dyadic depths; its Kraft sum and normalization bounds are
proved directly.

Separately, every selector which strictly prefers shorter words reverses some
preference whenever two unequal codeword lengths can be transposed.  This
transposition theorem has no finiteness premise.  A unary self-delimiting
codebook supplies a countably infinite witness.

The infinite boundary is explicit.  The unary codebook on $\mathbb N$ realizes
the geometric mass $2^{-(n+1)}$, whose total is one.  No strictly positive
uniform real mass is summable over $\mathbb N$, and no uncountable hypothesis
type injects into finite binary strings.  We do not claim a general countable
factor-two theorem without the corresponding infinite coding construction.

\subsection{Fixed vocabulary versus changing observer}

Bennett's continuous invariance theorem keeps the embodied vocabulary and
measure fixed and changes coordinates by a function-preserving
reparameterization.  If $F(\phi(\theta))=F(\theta)$, then the realized policy in
the fixed vocabulary is identical, so every extensional weakness readout is
identical.

This is compatible with observer relativity.  Holding the realized function
fixed while changing the vocabulary can change the visible policy cardinality.
The two results concern different transformations:
\[
\begin{array}{lll}
\text{coordinate change} & \text{fixed vocabulary and semantics}
  & \Longrightarrow \text{weakness invariant},\\
\text{observer change} & \text{different visible programs or distinctions}
  & \Longrightarrow \text{cardinality may change}.
\end{array}
\]
Measure-theoretic $\mu$-weakness and finite cardinal weakness are therefore not
identified.

\section{Layer-relative valuations}

Goertzel's quantale proposal treats weakness as a valuation-dependent algebraic
construction rather than a single scalar \cite{goertzel2026}.  The
formalization distinguishes exact recoveries from genuinely new profiles.

\subsection{Diagonal and complete-pair events}

The injective map $s\mapsto(s,s)$ defines a diagonal event.  Its cardinality is
exactly Bennett's admitted-world count, so its finite preference order is
preserved and reflected.  The complete-pair event $S\times S$ instead has
cardinality $|S|^2$.  A two-world canary gives diagonal cardinality two and
complete-pair cardinality four.

The diagonal here is the categorical diagonal in the category of types.  The
cardinality theorem is not a Lawvere diagonal argument: it has no evaluator,
self-application, weak point-surjectivity, fixed-point-free endomorphism, or
negation \cite{lawvere1969}.

\subsection{Logical entropy, graphtropy, and messages}

Logical entropy is treated as distinction probability following Ellerman
\cite{ellerman2017}.  Maximizing distinctions and minimizing distinctions are
opposite profiles over the same score.  On a Boolean carrier, discrete and
indiscrete partitions reverse preference between these profiles.

Graphtropy is defined over a declared distinction graph.  Adding distinction
edges decreases the missing-distinction score.  This fixes the edge polarity
in the type of the profile rather than hiding it in terminology.

For an MDL-like construction, a finite multiset is used as the free
commutative message monoid.  An atomic cost extends uniquely to an additive
message evaluator.  Unit atomic cost recovers ordinary message length and the
two-part model-plus-data comparison.  A nonuniform valuation reverses a length
comparison and does not factor through message cardinality.  This object is
called an additive valuation, not a quantale: completeness and distributivity
over arbitrary joins would require additional proofs.

\section{Epistemic resolve requires no bypass}

Crawford and Hammer describe epistemic resolve as a governed pathway from
recruitment through coupling, write-back, and downstream enforcement, together
with the absence of a parallel route that can change durable state outside the
pathway \cite{crawfordhammer2026}.  The last clause is not implied by checking
the local gates of one episode.

The formal model retains the information item, before and after states, and the
actual path of every update.  It separates:

\begin{itemize}
  \item local governance of an adjudicative episode;
  \item no bypass for any path producing the same durable transition or an
        equivalent downstream effect;
  \item strict resolve, in which every actual durable change has a resolving
        episode;
  \item architectural resolve, in which normal updates are governed and
        abnormal drift is detected and routed through a mediated correction.
\end{itemize}

\begin{theorem}[Gate passage is incomplete]
There is a finite architecture in which every local gate admits an update, yet
another path writes durable state and affects later behavior without coupling.
Hence local gate passage does not imply complete mediation.
\end{theorem}

A positive architecture has only the governed durable write and satisfies
strict resolve.  A corrective architecture satisfies the realistic
architectural condition while failing strict resolve, proving that the two
readings are genuinely different.

\section{What MHC says about GSLT structure}

The tempting analogy is to label one presentation ``systematic'', language
transport ``metasystematic'', and cross-language composition ``paradigmatic''.
The formal test rejects an unconditional assignment.

\subsection{Actual execution witnesses}

A binary operational schedule inside a GSLT contains a common initial term,
two state transformations, and four one-step witnesses for the two length-two
execution orders.  Endpoint semantics is a chain exactly when the endpoints
commute and a coordination exactly when they differ.

The interleaving product supplies a positive chain: a left component step and a
right component step form a commuting square and leave two simple children at
rank zero.  An authored interacting sum supplies a positive coordination:
alpha-then-beta and beta-then-alpha reach distinct terms and raise two simple
children to rank one.  Crucially, interaction was not used as a synonym for
coordination; endpoint inequality is proved.

\subsection{One naturality square, two MHC classifications}

An indexed operational GSLT retains two routes around its naturality square:
\[
\begin{aligned}
  r_1 &= \text{reduce in the source fibre, then transport},\\
  r_2 &= \text{transport, then reduce in the target fibre}.
\end{aligned}
\]
The routes have the same source and endpoint and are joined by a proof-relevant
two-cell.  Their first step constructors are nevertheless distinct.

\begin{theorem}[Observer-relative MHC classification]
For every indexed naturality square generated by a source step,
\begin{enumerate}[label=(\alph*)]
  \item endpoint observation classifies the two schedules as a chain;
  \item complete route observation classifies them as a coordination.
\end{enumerate}
\end{theorem}

The second result is not obtained from proof irrelevance: a finite first-step
observer separates an under-transport computation from an immediate transport
step, proving the complete routes unequal.

This is the tested divergence requested by the stack-alignment hypothesis.
GSLT presentation, transport, or cross-language composition alone determines
no context-free MHC stage.  A stage reading is licensed only after specifying
the outcome observer, proving actual schedule sensitivity or invariance, and
supplying homogeneous child ranks.  Endpoint equality and path identity answer
different questions, and neither should silently replace the other.

\section{Formal theorem map and scope}

\begin{center}
\small
\begin{tabular}{>{\raggedright\arraybackslash}p{.24\textwidth}
                >{\raggedright\arraybackslash}p{.31\textwidth}
                >{\raggedright\arraybackslash}p{.35\textwidth}}
\toprule
Module family & Positive result & Negative or boundary result \\
\midrule
Hierarchical complexity
  & arbitrary schedules, ordinal rank, finite recovery, $\varphi_n=2^n$
  & $\omega$ limit and atomless obstructions \\
Principal completion
  & product equivalence and finite cardinal product law
  & strict compatibility subfibre is not generally surjective \\
Enactive level bridge
  & unit-ranked paths give intrinsic levels, HC2 homogeneity, no gaps
  & transitive broad refinement admits shortcut lengths and no unit rank \\
Axis independence
  & complete-pair cardinality reduces to squared Bennett weakness
  & all four scalar axes fail pairwise factorization in both directions \\
Prior coding
  & finite factor-two construction and exact geometric infinite prior
  & recoding reversal, no positive uniform mass on $\mathbb N$, no uncountable codebook \\
Bennett reparameterization
  & fixed-vocabulary function-preserving invariance
  & changing the observer changes visible cardinality \\
Quantale razor valuations
  & diagonal recovery, explicit entropy and message profiles
  & polarity reversals and non-factorization under nonuniform valuation \\
Epistemic resolve
  & strict and corrective governed architectures
  & local gate passage admits a durable coupling bypass \\
GSLT hierarchical complexity
  & independent chain and noncommuting coordination witnesses
  & one naturality square changes class with its observer \\
\bottomrule
\end{tabular}
\end{center}

The development does not formalize empirical claims assigning human subjects
to developmental stages.  It formalizes the transferable coordination and
measurement structure.  It also does not claim that hierarchical order,
weakness, quantale value, description length, probability, or runtime cost are
universally reducible.  Each positive reduction above names the hypotheses that
make it true.

\section{Conclusion}

The common theory is not a new master scalar.  It is a discipline for keeping
informative structures separate until a theorem relates them.

MHC contributes an operational level-introduction rule: order rises when
permuting homogeneous lower actions changes the declared outcome.  Bennett
contributes semantic openness inside a fixed embodied language and a sharp
warning against treating code length as a representation-invariant ideal.
Goertzel contributes valuation-relative composition and exact specializations
that explain both agreement and reversal among weakness measures.  Crawford
and Hammer contribute the global no-bypass quantifier missing from local gate
checks.  GSLT execution contributes proof-relevant routes which reveal that
endpoint agreement does not erase computational organization.

Together these results support open-ended cognitive systems that retain their
paths, declare their observations, and earn each quotient or optimization by a
factorization, commutation, or mediation theorem.

\begin{thebibliography}{99}

\bibitem{commonsrichardskuhn1982}
M. L. Commons, F. A. Richards, and D. Kuhn.
Systematic and metasystematic reasoning: a case for a level of reasoning beyond
Piaget's stage of formal operations.
\emph{Child Development}, 53(4):1058--1069, 1982.

\bibitem{commonsetal1998}
M. L. Commons, E. J. Trudeau, S. A. Stein, F. A. Richards, and S. R. Krause.
Hierarchical complexity of tasks shows the existence of developmental stages.
\emph{Developmental Review}, 18:237--278, 1998.
doi:10.1006/drev.1998.0467.

\bibitem{commons2008}
M. L. Commons.
Introduction to the Model of Hierarchical Complexity and its relationship to
postformal action.
\emph{World Futures}, 64:305--320, 2008.
doi:10.1080/02604020802301105.

\bibitem{commonspekker2008}
M. L. Commons and A. Pekker.
Presenting the formal theory of hierarchical complexity.
\emph{World Futures}, 64:375--382, 2008.
doi:10.1080/02604020802301204.

\bibitem{bennett2026}
M. T. Bennett.
The Wrong Razor.
In M. Ikl\'e, A. Franz, M. Kemendo, and N. Lowy, editors,
\emph{Artificial General Intelligence}, LNAI 16854, pages 48--61, 2026.

\bibitem{goertzel2026}
B. Goertzel.
Weakness Is All You Need: Quantale Weakness as a Unifying Principle for
Cognition.
In M. Ikl\'e, A. Franz, M. Kemendo, and N. Lowy, editors,
\emph{Artificial General Intelligence}, LNAI 16854, pages 307--323, 2026.

\bibitem{crawfordhammer2026}
K. Crawford and P. Hammer.
Epistemic Resolve: Governed Uptake and the Architecture of Mind.
In M. Ikl\'e, A. Franz, M. Kemendo, and N. Lowy, editors,
\emph{Artificial General Intelligence}, LNAI 16854, pages 192--205, 2026.

\bibitem{ellerman2017}
D. Ellerman.
An introduction to logical entropy and its relation to Shannon entropy.
\emph{International Journal of Semantic Computing}, 11(2):121--145, 2017.

\bibitem{wallaceboulton1968}
C. S. Wallace and D. M. Boulton.
An information measure for classification.
\emph{The Computer Journal}, 11(2):185--194, 1968.

\bibitem{rissanen1978}
J. Rissanen.
Modeling by shortest data description.
\emph{Automatica}, 14(5):465--471, 1978.

\bibitem{lawvere1969}
F. W. Lawvere.
Diagonal arguments and cartesian closed categories.
In \emph{Category Theory, Homology Theory and Their Applications II},
Lecture Notes in Mathematics 92, pages 134--145. Springer, 1969.

\end{thebibliography}

\end{document}
